% MW-AeroFusion-Net: IEEE conference paper skeleton (PROGRESS P6.1). 6 pages.
% RULES (PROGRESS §0): no number in this file is typed by hand. Results come only from \input{tables/...} files that
% ml/scripts/make_tables.py generates from ml/runs/*/metrics.json (each row carries its run_id). Data facts below
% (station counts, periods) come from ml/data/stations_v1.csv, DATA_CARD.md and splits_v1.yaml.
% Wording: "learned attribution" / "source-group attribution gating"; never "causal". No "first"/"novel".
\documentclass[conference]{IEEEtran}
\usepackage{amsmath,amssymb,graphicx,booktabs,url,xcolor}
\newcommand{\todo}[1]{\textcolor{red}{[TODO: #1]}}

\begin{document}
\title{MW-AeroFusion-Net: A Leakage-Safe Multi-City Benchmark for Station-Level PM$_{2.5}$ Forecasting in India against CAMS}
\author{\IEEEauthorblockN{\todo{authors, group 40}}\IEEEauthorblockA{Dept. of CSE, BMS College of Engineering, Bengaluru, India}}
\maketitle

\begin{abstract}
\todo{write last, from the final tables only. Structure: problem (station-level PM$_{2.5}$ forecasts for Indian cities, 1--48 h);
gap (no leakage-safe benchmark against CAMS forecasts over CPCB stations, docs/related\_work.md); what we built (116-station
panel in 28 cities, three input settings, five-component model with component ablations and validated learned attribution);
results with run-traceable numbers; limitations.}
\end{abstract}

\section{Introduction}
Fine particulate matter (PM$_{2.5}$) is the dominant air-quality risk in Indian cities, and operational forecasts in India come
from chemistry-transport systems such as SAFAR~\cite{Beig_2021} and the IITM early-warning system~\cite{Ghude_2020,Jena_2021},
alongside global products such as the CAMS forecasts. Learned spatio-temporal models for PM$_{2.5}$ are well established
\cite{Wang_2020,Liang_2023,Yi_2018}, including graph models for Indian stations~\cite{Ejurothu_2022,Mandal_2023}, and
decomposition hybrids are common~\cite{Qiao_2019,Wang_2024}. Two practical questions remain open in the literature we
reviewed: whether such models beat a calibrated CAMS forecast at matched lead times under strictly causal inputs, and
which model components actually help.
\todo{contributions list, phrased as in docs/related\_work.md's novelty statement: (1) leakage-safe protocol and benchmark;
(2) MW-AeroFusion-Net with full component ladder and leave-one-out ablations; (3) learned attribution validated by
faithfulness, sanity and plausibility checks; (4) open code and station panel.}

\section{Related Work}
Graph models for traffic and air quality~\cite{li2017diffusion,Wu_2019,Yu_2018,Wang_2020}; attention models~\cite{Liang_2018,Liang_2023};
wavelet and decomposition hybrids and their leakage pitfall~\cite{Bai_2016,Qiao_2019,Quilty_2018,Barzegar_2021};
CAMS evaluation and post-processing~\cite{Wu_2020,Petetin_2022,Soni_2021}; interpretability~\cite{Lim_2021,adebayo2018sanity,jain2019attention}.
\todo{condense docs/related\_work.md tables A--E into $\sim$0.6 column; fill its `?' cells from full texts first.}

\section{Data}
\textbf{Targets.} Hourly CPCB PM$_{2.5}$ from the XKDR bulk export, converted from IST interval-start labels to UTC.
A fixed panel of 116 stations within 25\,km of 28 of the 33 study cities (stations reporting since 2020-01-01 with
$\geq$60\% hourly completeness); five cities without such a station are kept as unlabelled graph nodes. QC removes
sentinels, out-of-range values, flat-lines ($\geq$6\,h) and isolated spikes (trailing 24\,h Hampel rule, kept when
$\geq$2 stations in a city spike together). \textbf{Inputs.} ERA5 meteorology~\cite{Hersbach_2020} at nodes and on a
0.75$^\circ$ grid, FIRMS fire detections (counts, FRP, upwind counts), CAMS global forecasts, calendar and event flags.
\todo{data-card summary table generated from ml/data/DATA\_CARD.md.}

\section{Method}
\textbf{Graph.} Nodes are stations plus 17 upwind source regions; edges join the 8 nearest nodes within 600\,km with weight
$e^{-d/\sigma}\max(0,\cos(\theta_{\text{wind}}-\theta_{j\to i}))$ from past winds, plus a learned adaptive adjacency, and
diffusion convolution~\cite{li2017diffusion,Wu_2019}.
\textbf{Temporal encoder.} A causal Transformer~\cite{vaswani2017attention} over the past 72\,h per node (GRU in the ablation).
\textbf{Wavelet bands.} À-trous db4 bands computed causally inside each window (no full-series transform~\cite{Quilty_2018}).
\textbf{Regional context.} A ConvLSTM~\cite{shi2015convolutional} over gridded wind, boundary layer, FRP and CAMS fields,
sampled at node locations. \textbf{Source-group attribution gating.} A per-node softmax over feature groups (fire,
meteorology, dust, traffic proxy, CAMS, history, calendar) in the spirit of TFT variable selection~\cite{Lim_2021}, with an
entropy penalty. Its weights are a learned attribution, not a causal source apportionment.

\section{Experimental Protocol}
\textbf{Splits (UTC).} Train 2018--2021, validation 2022, test 2023--2024 (frozen, read once by a single script after model
freeze), a PM-only check April--September 2025, and a prospective winter 2026--27 test from live data.
\textbf{Latency rules.} Station data usable at $t_0-1$\,h; fire detections at acquisition$+3$\,h; a CAMS run $R$ from
$R+10$\,h. \textbf{Settings.} S-A past-only (primary); S-B adds the latest published CAMS forecast at matched leads;
S-C uses true future meteorology and is reported only as an upper bound. \textbf{Baselines.} Persistence, same hour
yesterday, climatology, raw CAMS, CAMS with per-station linear bias correction, LightGBM (S-A and MOS), LSTM, CNN-LSTM,
ConvLSTM, PM2.5-GNN with past-only meteorology~\cite{Wang_2020}, Graph WaveNet~\cite{Wu_2019}.
\textbf{Metrics.} RMSE, MAE, $R^2$, bias, skill against persistence and CAMS at +1/+24/+48\,h, by season, IGP vs rest,
and city; CSI for $>$250\,$\mu$g/m$^3$; AQI-class accuracy and macro-F1; Diebold--Mariano tests with HLN
correction~\cite{Diebold_1995} and a 24\,h block bootstrap; 5 seeds.

\section{Results}
\todo{\textbackslash input\{tables/test\_sa\_station.tex\} once final\_eval.py has run; until then only validation tables exist
(tables/val\_sa\_station.tex). Ladder table (M0--M6, val), ablation table (test, 5 seeds), season/IGP/episode tables.}

\section{Attribution Validation}
\todo{faithfulness (deletion vs random vs Integrated Gradients~\cite{sundararajan2017axiomatic}), cascading randomisation
\cite{adebayo2018sanity}, seed stability, plausibility (fire weight vs upwind FIRMS counts in Oct--Nov; dust weight vs CAMS
dust AOD in MAM), comparison with TreeSHAP on LightGBM~\cite{lundberg2017a}. Report failed checks as failed.}

\section{Limitations}
ERA5 stands in for operational analyses during training; the live system uses forecast meteorology. CPCB coverage after
2024 is sparse, so the 2025 check is PM-only. Five study cities have no long-running station. Attribution weights are
learned associations, validated but not causal. \todo{extend after results.}

\section{Conclusion}
\todo{from results only.}

\bibliographystyle{IEEEtran}
\bibliography{refs}
\end{document}
